% ============================================================
% 06_RESULTADOS.TEX
% Sistema Adaptativo de Classificação Multicritério
% de Fluidos de Corte
%
% FINALIDADE
%
% Estruturar a apresentação dos resultados que serão produzidos
% posteriormente pelos integrantes.
%
% Esta seção deverá receber, após a execução:
%
% 1. síntese da auditoria;
% 2. critérios finais;
% 3. pesos CRITIC;
% 4. ranking TOPSIS;
% 5. top 3;
% 6. conformidade;
% 7. cenários;
% 8. Pareto;
% 9. sensibilidade;
% 10. robustez.
%
% NENHUM resultado deve ser preenchido antecipadamente.
%
% ============================================================


% ============================================================
% FRAME 1 — VISÃO GERAL DOS RESULTADOS
% ============================================================

\begin{frame}{Resultados}

    \begin{block}{Resultados a serem apresentados}
        Esta seção será preenchida após a execução completa e validação
        da análise em R.
    \end{block}

    \vspace{0.3cm}

    A apresentação final deverá sintetizar:

    \begin{itemize}
        \item principais achados da auditoria;
        \item critérios utilizados;
        \item pesos obtidos pelo CRITIC;
        \item ranking obtido pelo TOPSIS;
        \item comportamento do top 3;
        \item conformidade técnica;
        \item análise de Pareto;
        \item sensibilidade;
        \item robustez.
    \end{itemize}

\end{frame}


% ============================================================
% FRAME 2 — RESULTADOS DA AUDITORIA
% ============================================================

\begin{frame}{Síntese da auditoria}

    \begin{block}{Principais resultados}
        \begin{itemize}
            \item Estrutura da base:
                  \texttt{TO\_BE\_FILLED}

            \item Inconsistências relevantes:
                  \texttt{TO\_BE\_FILLED}

            \item Reconciliações importantes:
                  \texttt{TO\_BE\_FILLED}

            \item Pendências:
                  \texttt{TO\_BE\_FILLED}

            \item Status da base canônica:
                  \texttt{TO\_BE\_FILLED}
        \end{itemize}
    \end{block}

    \vspace{0.2cm}

    \centering
    \textit{Apresentar apenas os achados que tenham impacto científico
    ou metodológico relevante.}

\end{frame}


% ============================================================
% FRAME 3 — CRITÉRIOS FINAIS
% ============================================================

\begin{frame}{Critérios utilizados}

    \begin{block}{Configuração final}
        \begin{itemize}
            \item Número de critérios:
                  \texttt{TO\_BE\_FILLED}

            \item Critérios incluídos:
                  \texttt{TO\_BE\_FILLED}

            \item Critérios excluídos relevantes:
                  \texttt{TO\_BE\_FILLED}

            \item Normalização principal:
                  \texttt{TO\_BE\_FILLED}
        \end{itemize}
    \end{block}

    \vspace{0.3cm}

    \begin{alertblock}{Interpretação}
        As exclusões deverão ser justificadas por critérios estatísticos,
        técnicos ou de qualidade dos dados, e não pelo efeito sobre o
        ranking.
    \end{alertblock}

\end{frame}


% ============================================================
% FRAME 4 — PESOS CRITIC
% ============================================================

\begin{frame}{Pesos obtidos pelo CRITIC}

    \begin{block}{Resultados}
        \begin{itemize}
            \item Critério de maior peso:
                  \texttt{TO\_BE\_FILLED}

            \item Critério de menor peso:
                  \texttt{TO\_BE\_FILLED}

            \item Soma dos pesos:
                  \texttt{TO\_BE\_FILLED}

            \item Concentração dos pesos:
                  \texttt{TO\_BE\_FILLED}
        \end{itemize}
    \end{block}

    \vspace{0.3cm}

    \begin{alertblock}{Cuidado}
        Peso elevado no CRITIC representa maior conteúdo informacional
        relativo no conjunto analisado, e não importância técnica
        absoluta.
    \end{alertblock}

\end{frame}


% ============================================================
% FRAME 5 — FIGURA FUTURA DOS PESOS
% ============================================================

\begin{frame}{Distribuição dos pesos CRITIC}

    \begin{center}

        \texttt{FIGURA\_CRITIC\_TO\_BE\_FILLED}

    \end{center}

    \vspace{0.3cm}

    \begin{block}{Interpretação}
        \texttt{TO\_BE\_FILLED}
    \end{block}

\end{frame}

% ------------------------------------------------------------
% Depois da execução, o bloco acima poderá ser substituído por:
%
% \begin{center}
%     \includegraphics[
%         width=0.90\textwidth
%     ]{figures/pesos_critic.pdf}
% \end{center}
%
% O nome real do arquivo deverá corresponder ao artefato
% efetivamente gerado pelos integrantes.
% ------------------------------------------------------------


% ============================================================
% FRAME 6 — RANKING TOPSIS
% ============================================================

\begin{frame}{Ranking multicritério}

    \begin{block}{Classificação principal}
        \begin{itemize}
            \item 1º lugar:
                  \texttt{TO\_BE\_FILLED}

            \item 2º lugar:
                  \texttt{TO\_BE\_FILLED}

            \item 3º lugar:
                  \texttt{TO\_BE\_FILLED}
        \end{itemize}
    \end{block}

    \vspace{0.3cm}

    Coeficientes de proximidade:

    \begin{center}
        \texttt{TO\_BE\_FILLED}
    \end{center}

    \vspace{0.2cm}

    \begin{alertblock}{Regra}
        O ranking oficial deverá ser obtido diretamente da execução
        aprovada do TOPSIS.
    \end{alertblock}

\end{frame}


% ============================================================
% FRAME 7 — FIGURA FUTURA DO RANKING
% ============================================================

\begin{frame}{Coeficientes TOPSIS}

    \begin{center}

        \texttt{FIGURA\_TOPSIS\_TO\_BE\_FILLED}

    \end{center}

    \vspace{0.3cm}

    \begin{block}{Mensagem principal}
        \texttt{TO\_BE\_FILLED}
    \end{block}

\end{frame}

% ------------------------------------------------------------
% Exemplo futuro:
%
% \includegraphics[
%     width=0.90\textwidth
% ]{figures/ranking_topsis.pdf}
%
% Não ativar enquanto o arquivo real não existir.
% ------------------------------------------------------------


% ============================================================
% FRAME 8 — GAP E PROXIMIDADE
% ============================================================

\begin{frame}{Diferenças entre as primeiras posições}

    \begin{block}{Top 3}
        \begin{itemize}
            \item Gap entre 1º e 2º:
                  \texttt{TO\_BE\_FILLED}

            \item Gap entre 2º e 3º:
                  \texttt{TO\_BE\_FILLED}

            \item Empates exatos:
                  \texttt{TO\_BE\_FILLED}

            \item Quase empates:
                  \texttt{TO\_BE\_FILLED}
        \end{itemize}
    \end{block}

    \vspace{0.3cm}

    \begin{block}{Interpretação}
        \texttt{TO\_BE\_FILLED}
    \end{block}

\end{frame}


% ============================================================
% FRAME 9 — PERFIL DO TOP 3
% ============================================================

\begin{frame}{Perfil das três primeiras alternativas}

    \begin{columns}[T]

        \begin{column}{0.31\textwidth}

            \begin{block}{1º lugar}
                \texttt{TO\_BE\_FILLED}

                \vspace{0.2cm}

                Pontos fortes:

                \texttt{TO\_BE\_FILLED}

                \vspace{0.2cm}

                Limitações:

                \texttt{TO\_BE\_FILLED}
            \end{block}

        \end{column}


        \begin{column}{0.31\textwidth}

            \begin{block}{2º lugar}
                \texttt{TO\_BE\_FILLED}

                \vspace{0.2cm}

                Pontos fortes:

                \texttt{TO\_BE\_FILLED}

                \vspace{0.2cm}

                Limitações:

                \texttt{TO\_BE\_FILLED}
            \end{block}

        \end{column}


        \begin{column}{0.31\textwidth}

            \begin{block}{3º lugar}
                \texttt{TO\_BE\_FILLED}

                \vspace{0.2cm}

                Pontos fortes:

                \texttt{TO\_BE\_FILLED}

                \vspace{0.2cm}

                Limitações:

                \texttt{TO\_BE\_FILLED}
            \end{block}

        \end{column}

    \end{columns}

\end{frame}


% ============================================================
% FRAME 10 — CONFORMIDADE
% ============================================================

\begin{frame}{Conformidade técnica}

    \begin{block}{Top 3}
        \begin{itemize}
            \item 1º lugar:
                  \texttt{TO\_BE\_FILLED}

            \item 2º lugar:
                  \texttt{TO\_BE\_FILLED}

            \item 3º lugar:
                  \texttt{TO\_BE\_FILLED}
        \end{itemize}
    \end{block}

    \vspace{0.3cm}

    \begin{block}{Restrições relevantes}
        \texttt{TO\_BE\_FILLED}
    \end{block}

    \vspace{0.2cm}

    \begin{alertblock}{Importante}
        Alto desempenho multicritério não elimina uma eventual
        não conformidade técnica.
    \end{alertblock}

\end{frame}


% ============================================================
% FRAME 11 — CENÁRIOS
% ============================================================

\begin{frame}{Resultados dos cenários}

    \begin{block}{Cenário principal}
        \texttt{TO\_BE\_FILLED}
    \end{block}

    \vspace{0.2cm}

    \begin{block}{Cenários alternativos}
        \texttt{TO\_BE\_FILLED}
    \end{block}

    \vspace{0.2cm}

    \begin{block}{Mudança do vencedor}
        \texttt{TO\_BE\_FILLED}
    \end{block}

    \vspace{0.2cm}

    \begin{block}{Estabilidade do top 3}
        \texttt{TO\_BE\_FILLED}
    \end{block}

\end{frame}


% ============================================================
% FRAME 12 — PARETO
% ============================================================

\begin{frame}{Análise de Pareto}

    \begin{block}{Alternativas não dominadas}
        \texttt{TO\_BE\_FILLED}
    \end{block}

    \vspace{0.3cm}

    \begin{block}{Status do primeiro colocado}
        \texttt{TO\_BE\_FILLED}
    \end{block}

    \vspace{0.3cm}

    \begin{block}{Interpretação}
        \texttt{TO\_BE\_FILLED}
    \end{block}

    \vspace{0.2cm}

    \begin{alertblock}{Cuidado}
        Uma alternativa não dominada não é automaticamente a melhor
        alternativa global.
    \end{alertblock}

\end{frame}


% ============================================================
% FRAME 13 — FIGURA FUTURA DE PARETO
% ============================================================

\begin{frame}{Fronteira de Pareto}

    \begin{center}

        \texttt{FIGURA\_PARETO\_TO\_BE\_FILLED}

    \end{center}

    \vspace{0.3cm}

    \begin{block}{Mensagem principal}
        \texttt{TO\_BE\_FILLED}
    \end{block}

\end{frame}


% ============================================================
% FRAME 14 — SENSIBILIDADE AOS PESOS
% ============================================================

\begin{frame}{Sensibilidade aos pesos}

    \begin{block}{Perturbações avaliadas}
        \texttt{TO\_BE\_FILLED}
    \end{block}

    \vspace{0.2cm}

    \begin{block}{Critério mais influente}
        \texttt{TO\_BE\_FILLED}
    \end{block}

    \vspace{0.2cm}

    \begin{block}{Mudança do primeiro colocado}
        \texttt{TO\_BE\_FILLED}
    \end{block}

    \vspace{0.2cm}

    \begin{block}{Mudança do top 3}
        \texttt{TO\_BE\_FILLED}
    \end{block}

\end{frame}


% ============================================================
% FRAME 15 — SENSIBILIDADE À NORMALIZAÇÃO
% ============================================================

\begin{frame}{Sensibilidade à normalização}

    \begin{block}{Métodos comparados}
        \texttt{TO\_BE\_FILLED}
    \end{block}

    \vspace{0.3cm}

    \begin{block}{Resultado principal}
        \texttt{TO\_BE\_FILLED}
    \end{block}

    \vspace{0.3cm}

    \begin{block}{Impacto sobre o ranking}
        \texttt{TO\_BE\_FILLED}
    \end{block}

\end{frame}


% ============================================================
% FRAME 16 — LEAVE-ONE-CRITERION-OUT
% ============================================================

\begin{frame}{Influência dos critérios}

    \begin{block}{Leave-one-criterion-out}
        \texttt{TO\_BE\_FILLED}
    \end{block}

    \vspace{0.3cm}

    \begin{block}{Critério cuja remoção mais alterou o ranking}
        \texttt{TO\_BE\_FILLED}
    \end{block}

    \vspace{0.3cm}

    \begin{block}{Interpretação}
        \texttt{TO\_BE\_FILLED}
    \end{block}

\end{frame}


% ============================================================
% FRAME 17 — LEAVE-ONE-ALTERNATIVE-OUT
% ============================================================

\begin{frame}{Influência do conjunto de alternativas}

    \begin{block}{Leave-one-alternative-out}
        \texttt{TO\_BE\_FILLED}
    \end{block}

    \vspace{0.3cm}

    \begin{block}{Rank reversal}
        \texttt{TO\_BE\_FILLED}
    \end{block}

    \vspace{0.3cm}

    \begin{block}{Alternativa de maior influência}
        \texttt{TO\_BE\_FILLED}
    \end{block}

\end{frame}


% ============================================================
% FRAME 18 — ROBUSTEZ
% ============================================================

\begin{frame}{Robustez do ranking}

    \begin{block}{Primeiro colocado}
        \texttt{TO\_BE\_FILLED}
    \end{block}

    \vspace{0.2cm}

    \begin{block}{Robustez local}
        \texttt{TO\_BE\_FILLED}
    \end{block}

    \vspace{0.2cm}

    \begin{block}{Robustez ampliada}
        \texttt{TO\_BE\_FILLED}
    \end{block}

    \vspace{0.2cm}

    \begin{block}{Estabilidade do top 3}
        \texttt{TO\_BE\_FILLED}
    \end{block}

\end{frame}


% ============================================================
% FRAME 19 — FIGURA FUTURA DE ROBUSTEZ
% ============================================================

\begin{frame}{Estabilidade das posições}

    \begin{center}

        \texttt{FIGURA\_ROBUSTEZ\_TO\_BE\_FILLED}

    \end{center}

    \vspace{0.3cm}

    \begin{block}{Mensagem principal}
        \texttt{TO\_BE\_FILLED}
    \end{block}

\end{frame}


% ============================================================
% FRAME 20 — FREQUÊNCIAS
% ============================================================

\begin{frame}{Frequência de posições nas configurações avaliadas}

    \begin{block}{Primeira posição}
        \texttt{TO\_BE\_FILLED}
    \end{block}

    \vspace{0.2cm}

    \begin{block}{Presença no top 3}
        \texttt{TO\_BE\_FILLED}
    \end{block}

    \vspace{0.3cm}

    \begin{alertblock}{Interpretação}
        Essas frequências representam o comportamento das alternativas
        nas configurações analisadas. Não são probabilidades reais de
        superioridade.
    \end{alertblock}

\end{frame}


% ============================================================
% FRAME 21 — SÍNTESE DOS RESULTADOS
% ============================================================

\begin{frame}{Síntese dos resultados}

    \begin{block}{Ranking principal}
        \texttt{TO\_BE\_FILLED}
    \end{block}

    \begin{block}{Conformidade}
        \texttt{TO\_BE\_FILLED}
    \end{block}

    \begin{block}{Pareto}
        \texttt{TO\_BE\_FILLED}
    \end{block}

    \begin{block}{Robustez}
        \texttt{TO\_BE\_FILLED}
    \end{block}

\end{frame}


% ============================================================
% FRAME 22 — MENSAGEM PARA A CONCLUSÃO
% ============================================================

\begin{frame}{O que os resultados permitem concluir?}

    Após todas as análises, deverá ser possível responder:

    \begin{itemize}
        \item qual alternativa ocupa a primeira posição?
        \item qual a diferença para as alternativas seguintes?
        \item a alternativa é tecnicamente conforme?
        \item ela pertence à fronteira de Pareto?
        \item a liderança permanece nos cenários?
        \item o ranking é sensível aos pesos?
        \item o ranking depende da normalização?
        \item há rank reversal?
        \item o top 3 é estável?
    \end{itemize}

    \vspace{0.3cm}

    \begin{block}{Próxima etapa}
        Integrar essas evidências em uma conclusão proporcional à
        robustez realmente observada.
    \end{block}

\end{frame}


% ============================================================
% REGRA DE SELEÇÃO PARA A APRESENTAÇÃO FINAL
% ============================================================
%
% Este arquivo contém mais frames do que provavelmente serão
% necessários na apresentação definitiva.
%
% Ele funciona como BASE.
%
% Após a execução, os integrantes deverão selecionar os slides
% que melhor respondem à narrativa científica.
%
% Não é necessário apresentar:
%
% - todos os cenários;
% - todas as perturbações;
% - todos os resultados LOCO;
% - todos os resultados LOAO;
% - todas as tabelas.
%
% Priorizar:
%
% resultado principal
% +
% evidência de interpretação
% +
% evidência de robustez
%
% ============================================================


% ============================================================
% FIGURAS E TABELAS
% ============================================================
%
% Resultados científicos deverão ser gerados pelos integrantes,
% preferencialmente de forma automatizada em R.
%
% Não digitar valores críticos diretamente no .tex quando houver
% mecanismo reproduzível para inseri-los.
%
% Pastas previstas no template da apresentação:
%
% templates/presentation/latex/figures/
%
% e, na apresentação final, poderão ser utilizadas as pastas
% correspondentes definidas pelo projeto.
%
% ============================================================


% ============================================================
% NÃO PREENCHER COM RESULTADOS PRELIMINARES
% ============================================================
%
% Resultados observados durante inspeções iniciais da planilha
% não devem ser tratados como resultados oficiais.
%
% Somente utilizar resultados associados a:
%
% base canônica aprovada
% +
% critérios aprovados
% +
% normalização aprovada
% +
% CRITIC validado
% +
% TOPSIS validado
% +
% robustez executada
%
% ============================================================


% ============================================================
% ESTADO ATUAL
% ============================================================
%
% RESULTS_SECTION_STATUS = BASE_PREPARED
%
% AUDIT_RESULTS = NOT_EXECUTED
% CRITIC_RESULTS = NOT_EXECUTED
% TOPSIS_RESULTS = NOT_EXECUTED
% SCENARIO_RESULTS = NOT_EXECUTED
% PARETO_RESULTS = NOT_EXECUTED
% SENSITIVITY_RESULTS = NOT_EXECUTED
% ROBUSTNESS_RESULTS = NOT_EXECUTED
%
% TOP_1 = TO_BE_FILLED
% TOP_2 = TO_BE_FILLED
% TOP_3 = TO_BE_FILLED
%
% ============================================================